\begin{theorem}[Degree three on every fixed nonordinary principal surface]
\label{thm:root-every-nonordinary-ppav-degree-three-full-candidate}
\status{Proved; complete consequence and independent dependency audits passed.}
Let $(A,\Theta)$ be a principally polarized abelian surface over a
perfect field of characteristic two, and assume its geometric
two-rank is zero or one. On the smooth maximal-rank open of the
complete linear system $|9\Theta|$, the relative unit local system
has rank eighty and both its geometric and arithmetic monodromy
groups are
\[
                            \GL_{80}(\mathbf F_2).
\]
Equivalently this is degree three for the embedding by $3\Theta$,
whenever the ambient hypersurface family surjects onto this complete
system, with the same restriction qualification as before.
\end{theorem}
\begin{proof}
We first work over an algebraic closure. If the principal surface
is a product, the previously proved complete supersingular-product
and mixed-product odd-degree theorems give exactly this group at
$d=3$. If it is nonproduct,
Lemma~\ref{lem:root-nonordinary-nonproduct-three-neighbor-candidate}
supplies a nonproduct nonordinary principal surface $B$ with a
polarized etale isogeny $B\to A$ of multiplier three. Apply
Theorem~\ref{thm:nonordinary-degree-three-from-neighbour-candidate}.
Its clean-boundary input is also checked independently in
Lemma~\ref{lem:root-nonordinary-three-theta-clean-boundary-candidate}.
The known all-degree principal-surface rank formula gives
$9\cdot3^2-1=80$, and the theorem supplies the full geometric
group in that dimension. Thus the conclusion holds for every
geometric principal surface, with no generic-modulus assumption.

The arithmetic image over the original perfect field contains
the geometric image in the same general linear group. Since that
geometric image is already the whole group, equality follows.
The chosen neighbour need only exist geometrically; no descent
of that auxiliary isogeny is required. If restriction from ambient
hypersurfaces to the complete system is surjective, its parameter
fibres are geometrically integral affine spaces on coefficient
charts, so the established restriction-transfer argument gives
the same monodromy there.
\end{proof}
